A non-divergem form of consevation equations  can be obtained 
using Howarth-Dorodnitsyn transformation. 
Definig the stream function for compressible flow as: 

\begin{equation}
\frac{\partial \Psi}{\partial x}=-\rho u \qquad
\frac{\partial \Psi}{\partial y}=\rho v,
\end{equation} 

The mass conservation equation is identically satisfied by 
the stream function.
%
%
\begin{equation}
\frac{\partial{(\rho u)}}{\partial x}
+
\frac{\partial{(\rho v)}}{\partial y}
=0
\end{equation}

\begin{equation}
-\frac{\partial^2 \Psi }{\partial  \overline{x}\partial \overline{y}}
+
\frac{\partial^2 \Psi }{\partial \overline{x}\partial \overline{y}}
=0
\end{equation}

Now defining a vertical stretching as:

\begin{equation}
\epsilon =x\qquad
\eta=\int_0^{y}\rho dy%\qquad 
%\hat{V}=\varrho V+U\int_o^{\overline{y}}\varrho_{\overline{x}}\qquad
\end{equation} 

and using the chain rule, a function

\begin{equation}
f
=
f(\epsilon,\eta)=
f(\epsilon(x,y),\eta(x,y))
\end{equation} 


and its derivatives 
can be expressed by

\begin{equation}
\frac{\partial f}{\partial x}
=
\frac{\partial f}{\partial \epsilon}
\frac{\partial \epsilon}{\partial x}
+
\frac{\partial f}{\partial \eta}
\frac{\partial \eta}{\partial x}
\end{equation} 

and 

\begin{equation}
\frac{\partial f}{\partial y}=
\frac{\partial f}{\partial \epsilon}
\frac{\partial \epsilon}{\partial y}
+
\frac{\partial f}{\partial \eta}
\frac{\partial \eta}{\partial y}
.
\end{equation} 

Then,

\begin{equation}
\frac{\partial f}{\partial x}=
\frac{\partial f}{\partial \epsilon}
+
\frac{\partial f}{\partial \eta}
\int_0^{y}\frac{\partial \rho }{\partial x}dy
,
\end{equation} 

and 

\begin{equation}
\frac{\partial f}{\partial y}
=
\rho
\frac{\partial f}{\partial \eta}
.
\end{equation} 

Using the stream function definition and the 
H-D transformation  must be defined a new velocity $U$ as:

\begin{equation}
\frac{\partial \Psi}{\partial y}
=
\rho u
\end{equation} 

\begin{equation}
\rho
\frac{\partial \Psi}{\partial \eta}
=
\rho
u
\end{equation} 

\begin{equation}
\frac{\partial \Psi}{\partial \eta}
=
u
=
U
\end{equation} 

that have the same form as incompresible flow. 
In the same way was defined

\begin{equation}
\rho v=-\frac{\partial \Psi}{\partial x}.
\end{equation} 

With  help of the transformation equations for  $(x,y)$  to 
$(\eta,\epsilon)$ : 

\begin{equation}
  \rho v=
  -
  \left(
       \frac{\partial \Psi}{\partial \epsilon}+
       \frac{\partial \Psi}{\partial \eta}
       \frac{\partial \eta}{\partial x}
  \right),
\end{equation} 

then 
\begin{equation}
  \rho v=
       V
       -
       U
       \frac{\partial \eta}{\partial x}
\end{equation} 

Defining 

\begin{equation}
       V=-\frac{\partial \Psi}{\partial \epsilon}
       ,
\end{equation} 

then 

\begin{equation}
    V
    =
    \rho v
    +
    U
    \frac{\partial \eta}{\partial x}
    =
    \rho v
    +
    U
    \int_0^y
    \frac{\partial \rho}{\partial x}
    .
\end{equation} 

Using  the variable change and its 
consequent new velocity definitions, the conservation 
equation becomes: 

%%%\begin{equation} 
%%%\frac{\partial (\rho u)}{\partial x}
%%%+
%%%\frac{\partial (\rho v)}{\partial y}
%%%=
%%%0
%%%\end{equation} 
%%%
%%%
%%%\begin{equation} 
%%%\frac{\partial (\rho u)}{\partial \epsilon}
%%%+
%%%\frac{\partial (\rho u)}{\partial \eta}
%%%\frac{\partial \eta}{\partial x }
%%%+
%%%\rho
%%%\frac{\partial (\rho v)}{\partial \eta}
%%%=
%%%0
%%%\end{equation} 
%%%
%%%\begin{equation} 
%%%\frac{\partial (\rho U)}{\partial \epsilon}
%%%+
%%%\frac{\partial (\rho U)}{\partial \eta}
%%%\frac{\partial \eta}{\partial x }
%%%+
%%%\rho
%%%\frac{\partial }{\partial \eta}
%%%\left(
%%%    V
%%%    -
%%%    U
%%%    \frac{\partial \eta}{\partial x }
%%%\right)
%%%=
%%%0
%%%\end{equation} 
%%%
%%%\begin{equation} 
%%%\rho 
%%%\frac{\partial U}{\partial \epsilon}
%%%+
%%%U
%%%\frac{\partial \rho }{\partial \epsilon}
%%%+
%%%\rho
%%%\frac{\partial  U}{\partial \eta}
%%%\frac{\partial \eta}{\partial x }
%%%+
%%%U
%%%\frac{\partial  \rho}{\partial \eta}
%%%\frac{\partial \eta}{\partial x }
%%%+
%%%\rho
%%%\frac{\partial V}{\partial \eta}
%%%-
%%%\rho
%%%\frac{\partial}{\partial \eta}
%%%\left(
%%%    U
%%%    \frac{\partial \eta}{\partial x }
%%%\right)
%%%=
%%%0
%%%\end{equation} 
%%%
%%%\begin{equation} 
%%%\rho 
%%%\frac{\partial U}{\partial \epsilon}
%%%+
%%%U
%%%\frac{\partial \rho }{\partial \epsilon}
%%%+
%%%\cancel{
%%%\rho
%%%\frac{\partial  U}{\partial \eta}
%%%\frac{\partial \eta}{\partial x }
%%%}
%%%+
%%%U
%%%\frac{\partial  \rho}{\partial \eta}
%%%\frac{\partial \eta}{\partial x }
%%%+
%%%\rho
%%%\frac{\partial V}{\partial \eta}
%%%-
%%%\cancel{
%%%\rho
%%%\frac{\partial \eta}{\partial x }
%%%\frac{\partial}{\partial \eta}
%%%\left(
%%%    U
%%%\right)
%%%}
%%%-
%%%\rho
%%%U
%%%\frac{\partial}{\partial \eta}
%%%\left(
%%%    \frac{\partial \eta}{\partial x }
%%%\right)
%%%=
%%%0
%%%\end{equation} 
%%%
%%%
%%%\begin{equation} 
%%%\rho 
%%%\frac{\partial U}{\partial \epsilon}
%%%+
%%%U
%%%\frac{\partial \rho }{\partial \epsilon}
%%%+
%%%U
%%%\frac{\partial  \rho}{\partial \eta}
%%%\frac{\partial \eta}{\partial x }
%%%+
%%%\rho
%%%\frac{\partial V}{\partial \eta}
%%%-
%%%U
%%%\frac{\partial}{\partial y}
%%%\left(
%%%    \frac{\partial \eta}{\partial x }
%%%\right)
%%%=
%%%0
%%%\end{equation} 
%%%
%%%Using the fact that 
%%%
%%%\begin{equation*} 
%%%\frac{\partial}{\partial y}
%%%\left(
%%%    \frac{\partial \eta}{\partial x }
%%%\right)
%%%=
%%%\frac{\partial}{\partial y}
%%%    \int_0^y
%%%    \frac{\partial \rho}{\partial x}
%%%    dy
%%%    =
%%%    \frac{\partial \rho}{\partial x}
%%%    =
%%%    \frac{\partial \rho}{\partial \epsilon}
%%%    +
%%%    \frac{\partial \rho}{\partial \eta}
%%%    \frac{\partial \eta}{\partial x}
%%%\end{equation*} 
%%%
%%%then 
%%%
%%%\begin{equation} 
%%%\rho 
%%%\frac{\partial U}{\partial \epsilon}
%%%+
%%%U
%%%\frac{\partial \rho }{\partial \epsilon}
%%%+
%%%\cancel{
%%%U
%%%\frac{\partial  \rho}{\partial \eta}
%%%\frac{\partial \eta}{\partial x }
%%%}
%%%+
%%%\cancel{
%%%\rho
%%%\frac{\partial V}{\partial \eta}
%%%}
%%%-
%%%\cancel{
%%%U
%%%\frac{\partial}{\partial y}
%%%\left(
%%%    \frac{\partial \eta}{\partial x }
%%%\right)
%%%}
%%%=
%%%0
%%%\end{equation} 
%%%
%%%The mass conservation equation for a compressible boundary 
%%%layer can be written is a non-divergent form as:
%%%
%%%\begin{equation} 
%%%\rho 
%%%\frac{\partial U}{\partial \epsilon}
%%%+
%%%\rho
%%%\frac{\partial V}{\partial \eta}
%%%=
%%%0
%%%\end{equation} 
%%%
%%%xx

\begin{equation} 
\frac{\partial (\rho u)}{\partial x}
+
\frac{\partial (\rho v)}{\partial y}
=
0
\end{equation} 

\begin{equation} 
\frac{\partial (\rho U)}{\partial x}
+
\frac{\partial }{\partial y}
\left(
    V
    -
    U
    \frac{\partial \eta}{\partial x }
\right)
=
0
\end{equation} 

\begin{equation} 
\frac{\partial (\rho U)}{\partial x}
+
\frac{\partial V}{\partial y}
-
\frac{\partial }{\partial y}
\left(
    U
    \frac{\partial \eta}{\partial x }
\right)
=
0
\end{equation} 

\begin{equation} 
\rho
\frac{\partial  U}{\partial x}
+
\frac{\partial V}{\partial y}
+
U
\frac{\partial \rho }{\partial x}
-
U
\frac{\partial }{\partial y}
\left(
    \frac{\partial \eta}{\partial x }
\right)
-
\frac{\partial \eta}{\partial x }
\frac{\partial U}{\partial y}
=
0
\end{equation} 

Using the fact 

\begin{equation*} 
\frac{\partial}{\partial y}
\left(
    \frac{\partial \eta}{\partial x }
\right)
=
\frac{\partial}{\partial y}
    \int_0^y
    \frac{\partial \rho}{\partial x}
    dy
    =
    \frac{\partial \rho}{\partial x}
\end{equation*} 

\begin{equation} 
\rho
\frac{\partial  U}{\partial x}
+
\frac{\partial V}{\partial y}
-
\frac{\partial \eta}{\partial x }
\frac{\partial U}{\partial y}
=
0
\end{equation} 

Now,  changing the variables 

\begin{equation} 
\rho
\frac{\partial  U}{\partial \epsilon}
+
\frac{\partial  U}{\partial \eta}
\frac{\partial  \eta}{\partial x}
+
\rho
\frac{\partial V}{\partial \eta}
-
\frac{\partial \eta}{\partial x }
\frac{\partial U}{\partial y}
=
0
.
\end{equation} 

The mass conservation equation for a compressible boundary 
layer can be written is a non-divergent form as:

\begin{equation} 
\rho 
\frac{\partial U}{\partial \epsilon}
+
\rho
\frac{\partial V}{\partial \eta}
=
0
.
\end{equation} 

For the consevation species, a similar process 
allows to obtain: 

\begin{equation} 
\rho u\frac{\partial Y_{i}}{\partial \epsilon}
+
\rho u\frac{\partial Y_{i}}{\partial \eta}
\frac{\partial \eta}{\partial  x}
+
v\frac{\partial Y_{i}}{\partial \eta}
=
\frac{\partial}{\partial y}
\left(\mu \frac{\partial Y_i}{\partial y}\right)
-
\dot{\omega}_i
,
\end{equation} 

\begin{equation} 
\rho u\frac{\partial Y_{i}}{\partial \epsilon}
+
\rho u\frac{\partial Y_{i}}{\partial \eta}
\frac{\partial \eta}{\partial  x}
+
\rho v
\frac{\partial Y_{i}}{\partial y}
=
\frac{\partial}{\partial y}
\left(\mu \frac{\partial Y_i}{\partial y}\right)
-
\dot{\omega}_i
,
\end{equation} 

\begin{equation} 
U
\frac{\partial Y_{i}}{\partial \epsilon}
+
V
\frac{\partial Y_{i}}{\partial \eta}
=
\frac{\partial}{\partial \eta}
\left(\mu\rho \frac{\partial Y_i}{\partial \eta}\right)
-
\frac{\dot{\omega}_i}{\rho}
,
\end{equation} 

In the same way for the streamwise direction: 

\begin{equation} 
U\frac{\partial U}{\partial \epsilon}
+
V\frac{\partial U}{\partial \eta}
=-
\frac{\partial}{\partial \eta}
\left(\mu\rho\frac{\partial u}{\partial \eta}\right)
,
\end{equation} 

and the total enthalpy equation:

\begin{equation*} 
\begin{split} 
U 
\frac{\partial H}{\partial \epsilon}
+
V 
\frac{\partial H}{\partial \eta}
=
\frac{\partial}{\partial \eta}
    \left[
    \mu
    \rho
    \frac{\partial H}{\partial \eta}
    \right]
    ,
\end{split} 
\end{equation*} 

and 

\begin{equation} 
%p^*=\rho^*R^*T^*
%p\rho_1*u_1^2=\rho^*_1 R^* T_1^*\rho T
%p u_1^2=\frac{\gamma}{\gamma}R^*T_1^*\rho T
p=\frac{1}{\gamma M^2} \rho T
.
\end{equation} 

In function of the stream function for compressible flow

\begin{equation} 
U
=
\frac{\partial \psi}{\partial \eta}
,
\qquad
V
=
-\frac{\partial \psi}{\partial \epsilon}
\end{equation} 

%\newpage

the consevation equations become
:

\begin{equation} 
\frac{\partial \psi}{\partial \eta}
\frac{\partial Y_{i}}{\partial \epsilon}
-
\frac{\partial \psi}{\partial \epsilon}
\frac{\partial Y_{i}}{\partial \eta}
=
\frac{\partial}{\partial \eta}
\left(\mu\rho \frac{\partial Y_i}{\partial \eta}\right)
-
\frac{\dot{\omega}_i}{\rho}
,
\end{equation} 


\begin{equation} 
\frac{\partial \psi}{\partial \eta}
\frac{\partial U}{\partial \epsilon}
-
\frac{\partial \psi}{\partial \epsilon}
\frac{\partial U}{\partial y}
=-
\frac{\partial}{\partial y}
\left(\mu\rho\frac{\partial u}{\partial \eta}\right)
,
\end{equation} 


\begin{equation*} 
\begin{split} 
\frac{\partial \psi}{\partial \eta}
\frac{\partial H}{\partial \epsilon}
-
\frac{\partial \psi}{\partial \epsilon}
\frac{\partial H}{\partial \epsilon}
=
\frac{\partial}{\partial \eta}
    \left(
    \mu
    \rho
    \frac{\partial H}{\partial \eta}
    \right)
    ,
\end{split} 
\end{equation*} 

The species equation can be writtem in terms  of 
passive scalar  
$Z=(Z_1,Z_2,Z_3)$. 
Zeldovich proposed using the $Z$ variable, 
which is a linear combination of the
mass fractions of the fuel and oxidizer, to eliminnate 
the reaction rate . 
Consider a single step chemical reaction of the type $\nu_F F+\nu_OO \rightarrow \nu_P P$
where  

\begin{equation}
\dot{\omega}_F=-W_F\nu_F\dot{\Omega}
,
\quad              
\dot{\omega}_O=-W_O\nu_O\dot{\Omega}
,
\quad              
\dot{\omega}_P= W_P\nu_P\dot{\Omega}
.
\end{equation}

Normalized the coeficcient by $\nu_F$: 
$F+\dfrac{\nu_O}{\nu_F}O \rightarrow \dfrac{\nu_P}{\nu_F} P$. 

To transform into mass form multiplie by 
their respectively molecular way, once  $m_i=W_i\nu_i$,   
$F+\dfrac{W_O\nu_O}{W_F\nu_F}O \rightarrow \dfrac{W_P\nu_P}{W_F\nu_F} P$. 

Now  

\begin{equation}
    \frac{\dot{\omega_F}}{\dot{\omega_O}}
    =
    \frac{W_F\nu_F\dot{\Omega}}{W_O\nu_O\dot{\Omega}}
    =
    \frac{1}{n} 
\end{equation}

with $n=(W_O\nu_O)/W_F\nu_F$.
So,  
\begin{equation}
    \dot{\omega_F}=\frac{1}{n}{\dot{\omega_O}}
\end{equation}

Then, the single step chemical reaction gets  $F + n O \rightarrow (n+1)P$. 
Definning the passive scalars

\begin{equation}
Z_1=
Y_F
-
\frac{Y_O}{n}
\end{equation}          

which is a linear combination of the fuel and oxidizer mass fractions, 
meaning that for stiochiometric conditions, 
measures how far the mixture is from stoichiometric balance.
if $Z_1=0$ the mixture is stoichiometric, for a 
fueled mixture, $Z_1>1$, and for an oxidizer rich mixture, $Z_1<1$, 
Then:

\begin{equation} 
\frac{\partial \psi}{\partial \eta}
\frac{\partial Y_{F}}{\partial \epsilon}
-
\frac{\partial \psi}{\partial \epsilon}
\frac{\partial Y_{F}}{\partial \eta}
-
\frac{1}{n}
\frac{\partial \psi}{\partial \eta}
\frac{\partial Y_{O}}{\partial \epsilon}
+
\frac{\partial \psi}{\partial \epsilon}
\frac{\partial Y_{O}}{\partial \eta}
=
\frac{\partial}{\partial \eta}
\left(\mu\rho \frac{\partial Y_F}{\partial \eta}\right)
-
\frac{\partial}{\partial \eta}
\left(\mu\rho \frac{\partial Y_O}{\partial \eta}\right)
-
\frac{\dot{\omega}_F}{\rho}
+
\frac{\dot{\omega}_O}{n\rho}
.
\end{equation} 

\begin{equation} 
\frac{\partial \psi}{\partial \eta}
\frac{\partial Z_1}{\partial \epsilon}
-
\frac{\partial \psi}{\partial \epsilon}
\frac{\partial Z_1}{\partial \eta}
=
\frac{\partial}{\partial \eta}
\left(\mu\rho \frac{\partial Z_1}{\partial \eta}\right)
-
\frac{1}{\rho}
\left(
    \dot{\omega}_F
    -
    \frac{\dot{\omega}_O}{n}
\right)
,
\end{equation} 

with 

\begin{equation}
    \dot{\omega_F}=\frac{1}{n}{\dot{\omega_O}}
    ,
\end{equation}

the $Z_1$ equation gets 

\begin{equation} 
\frac{\partial \psi}{\partial \eta}
\frac{\partial Z_1}{\partial \epsilon}
-
\frac{\partial \psi}{\partial \epsilon}
\frac{\partial Z_1}{\partial \eta}
=
\frac{\partial}{\partial \eta}
\left(
    \mu\rho  \frac{\partial Z_1}{\partial \eta}
\right)
.
\end{equation} 

Similarly, another passive scalar can be defined as

\begin{equation}
%Z=\frac{Y_F-Y_O+Y_{O\infty}}{Y_{F\infty}+Y_{O\infty}}
Z_2=
Y_F
+
\frac{Y_P}{n+1}
\end{equation}          

which is a linear combination of the fuel and product mass fractions,
measuring how far the mixture is from the pure fuel. 
$Z_2=0$ only for a pure oxidizer flow. 
The consevation equation for $Z_2$ has the same form that $Z_1$:  

\begin{equation} 
\frac{\partial \psi}{\partial \eta}
\frac{\partial Z_2}{\partial \epsilon}
-
\frac{\partial \psi}{\partial \epsilon}
\frac{\partial Z_2}{\partial \eta}
=
\frac{\partial}{\partial \eta}
\left(
     \mu\rho  \frac{\partial Z_2}{\partial \eta}
\right)
.
\end{equation} 

This is due that: $F+n O \rightarrow (n+1)P$

\begin{equation}
    \frac{\dot{\omega_F}}{\dot{\omega_P}}
    =
    \frac{W_F\nu_F\dot{\Omega}}{W_P\nu_P\dot{\Omega}}
    =
    \frac{1}{(n+1)}
\end{equation}

with $(n+1)=(W_P\nu_P)/W_F\nu_F$ 

So,  
\begin{equation}
    \dot{\omega_F}=\frac{1}{(n+1)}{\dot{\omega_P}}
\end{equation}
 
For the oxidizer, another passive scalar can be defined as

\begin{equation}
Z_3=
\frac{Y_O}{n}
+
\frac{Y_P}{n+1}
\end{equation}          

which is a linear combination of the oxidizer and product mass fractions,
measuring how far the mixture is from the pure oxidizer. 
$Z_3=0$ only for a pure fuel flow. 

In the same way: 

\begin{equation}
    \frac{\dot{\omega_O}}{\dot{\omega_P}}
    =
    \frac{W_O\nu_O\dot{\Omega}}{W_P\nu_P\dot{\Omega}}
    =
    \frac{n}{(n+1)} 
\end{equation}

with $n_p=n/n+1=W_P\nu_P/(W_O\nu_O)$ 

So,  
\begin{equation}
    \dot{\omega_P}=\frac{1}{n_p}{\dot{\omega_O}}
    =\frac{n+1}{n}{\dot{\omega_O}}
\end{equation}

The consevation equation hold the same form as the previouly defined: 

\begin{equation} 
\frac{\partial \psi}{\partial \eta}
\frac{\partial Z_3}{\partial \epsilon}
-
\frac{\partial \psi}{\partial \epsilon}
\frac{\partial Z_3}{\partial \eta}
=
\frac{\partial}{\partial \eta}
\left(
     \mu\rho  \frac{\partial Z_3}{\partial \eta}
\right)
.
\end{equation} 

\begin{equation}
Pr
=
\frac{\mu c_p}{k}
=\frac{\mu \rho c_p}{k\rho}
\end{equation} 


The complete system of equations for a compressible 
reacting mixing layer is: 

\begin{equation} 
\frac{\partial \psi}{\partial \eta}
\frac{\partial Z_{i}}{\partial \epsilon}
-
\frac{\partial \psi}{\partial \epsilon}
\frac{\partial Z_{i}}{\partial \eta}
=
\frac{\partial}{\partial \eta}
\left(\mu\rho  \frac{\partial Z_i}{\partial \eta}\right)
\end{equation} 


\begin{equation} 
\frac{\partial \psi}{\partial \eta}
\frac{\partial U}{\partial \epsilon}
-
\frac{\partial \psi}{\partial \epsilon}
\frac{\partial U}{\partial y}
=-
\frac{\partial}{\partial y}
\left(\mu\rho\frac{\partial u}{\partial \eta}\right)
,
\end{equation} 


\begin{equation*} 
\begin{split} 
\frac{\partial \psi}{\partial \eta}
\frac{\partial H}{\partial \epsilon}
-
\frac{\partial \psi}{\partial \epsilon}
\frac{\partial H}{\partial \epsilon}
=
\frac{\partial}{\partial \eta}
    \left(
    \mu
    \rho
        \frac{\partial H}{\partial \eta}
    \right)
    ,
\end{split} 
\end{equation*} 

The gas ideal equation are not necessary,  
the pressure is decople for the system. 

Other change of variable  can be applied
in orther to simplified the system, 
defining: 

\begin{equation*} 
    (\epsilon,\eta)
    \rightarrow
    (\chi,\gamma)
    ,
    \qquad
    \text{with}
    \qquad
    \chi=\epsilon
    ,
    \qquad
    \gamma
    =
    \frac{\eta}{\sqrt{2\epsilon}}
    ,
    \qquad
    \epsilon=\epsilon(\chi,\gamma), 
    \qquad
    \eta=\eta(\chi,\gamma), 
\end{equation*} 

Now 

\begin{equation*}
    f=f(\epsilon(\chi,\gamma),\eta(\chi,\gamma)=f(\chi,\gamma))
\end{equation*} 

and using the chain rule for  the derivates of the function $f$, 

\begin{equation}
\left(
\frac{\partial f}{\partial \epsilon}
\right)_\eta
=
\left(
\frac{\partial f}{\partial \chi}
\right)_\gamma
\left(
\frac{\partial \chi}{\partial \epsilon}
\right)_\eta
+
\left(
\frac{\partial f}{\partial \gamma}
\right)_\chi
\left(
\frac{\partial \gamma}{\partial \epsilon}
\right)_\eta
=
\left(
\frac{\partial f}{\partial \chi}
\right)_\gamma
+
\left(
\frac{\partial f}{\partial \gamma}
\right)_\chi
\left(
\frac{\partial \gamma}{\partial \epsilon}
\right)_\eta
\end{equation} 

with

\begin{equation}
\left(
\frac{\partial \gamma}{\partial \epsilon}
\right)_\eta
=
\frac{\partial }{\partial \chi}
\left(
    \frac{\eta}{\sqrt{2\epsilon}}
\right)
=
\eta
\frac{\partial }{\partial \chi}
\left(
    (2\epsilon)^{-1/2}
\right)
=
-
\eta
(2\epsilon)^{-3/2}
=
-
\frac{\eta}{(2\epsilon)^{1/2}(2\epsilon)}
=
-
\frac{\gamma}{2\epsilon}
\end{equation} 

then 

\begin{equation}
\left(
\frac{\partial f}{\partial \epsilon}
\right)_\eta
=
\left(
\frac{\partial f}{\partial \chi}
\right)_\gamma
-
\frac{\gamma}{2\epsilon}
\left(
\frac{\partial f}{\partial \gamma}
\right)_\chi
\end{equation} 

and 

\begin{equation}
\left(
\frac{\partial f}{\partial \eta}
\right)_\epsilon
=
\left(
\frac{\partial f}{\partial \chi}
\right)_\gamma
\left(
\frac{\partial \chi}{\partial \eta}
\right)_\epsilon
+
\left(
\frac{\partial f}{\partial \gamma}
\right)_\chi
\left(
\frac{\partial \gamma}{\partial \eta}
\right)_\epsilon
=
\left(
\frac{\partial f}{\partial \gamma}
\right)_\chi
\frac{1}{\sqrt{2\epsilon}}
.
\end{equation} 

A function $F$ is assumed as:

\begin{equation} 
F(\chi,\gamma)=\frac{\psi(\epsilon,\eta)}{\sqrt{2\epsilon}}
\end{equation} 

Now using in the momentum streamwise direction with 

\begin{equation*} 
\frac{\partial \psi}{\partial \eta}
\frac{\partial U}{\partial \epsilon}
-
\frac{\partial \psi}{\partial \epsilon}
\frac{\partial U}{\partial \eta}
=
\frac{\partial}{\partial \eta}
\left(\mu\rho\frac{\partial U}{\partial \eta}\right)
,
\end{equation*} 

\begin{equation*} 
\frac{\partial \psi}{\partial \eta}
\left(
    \frac{\partial U}{\partial \chi}
    -
    \frac{\gamma}{2\epsilon}
    \frac{\partial U}{\partial \gamma}
\right)
-
\frac{\partial \psi}{\partial \epsilon}
\left(
    \frac{1}{\sqrt{2\epsilon}}
    \frac{\partial U}{\partial \gamma}
\right)
=
\frac{1}{\sqrt{2\epsilon}}
\frac{\partial}{\partial \gamma}
\left[
    \mu\rho
    \left(
    \frac{1}{\sqrt{2\epsilon}}
    \frac{\partial U}{\partial \gamma}
    \right)
\right]
,
\end{equation*} 

The first and second derivate  of  $F$ are given by 

\begin{equation*} 
F'=
\frac{\partial F(\chi,\gamma)}{\partial \gamma}
=
\frac{\partial }{\partial \eta}
\left(
    \frac{\psi(\epsilon,\eta)}{\sqrt{2\epsilon}}
\right)
\frac{\partial \eta}{\partial \gamma}
=
\frac{\partial }{\partial \eta}
\left(
    \psi(\epsilon,\eta)
\right)
\frac{\sqrt{2\epsilon}}
{\sqrt{2\epsilon}}
=
U
\end{equation*}  

\begin{equation*} 
F^{\prime\prime}=
\frac{\partial }{\partial \gamma}
\left(
\frac{\partial F(\chi,\gamma)}{\partial \gamma}
\right)
=
\frac{\partial U}{\partial \gamma}
\end{equation*} 

The derivate of $F$ respect to $\chi$

\begin{equation*} 
\begin{split}  
\left(
    \frac{\partial F(\chi,\gamma)}{\partial \chi}
\right)_\gamma
=
\left(
\frac{\partial F(\chi,\gamma)}{\partial \epsilon}
\right)_\eta
\left(
\frac{\partial \epsilon}{\partial \chi}
\right)_\gamma
+
\left(
\frac{\partial F(\chi,\gamma)}{\partial \eta}
\right)_\epsilon
\left(
\frac{\partial \eta}{\partial \chi}
\right)_\gamma
\\
=
\left(
\frac{\partial F(\chi,\gamma)}{\partial \epsilon}
\right)_\eta
+
\left(
\frac{\partial F(\chi,\gamma)}{\partial \eta}
\right)_\epsilon
\left(
\frac{\partial \eta}{\partial \epsilon}
\right)_\gamma
\left(
\frac{\partial \epsilon}{\partial \chi}
\right)_\gamma
=
\\
\left(
\frac{\partial F(\chi,\gamma)}{\partial \epsilon}
\right)_\eta
+
\left(
\frac{\partial F(\chi,\gamma)}{\partial \eta}
\right)_\epsilon
\left(
\frac{\partial \eta}{\partial \epsilon}
\right)_\gamma
\end{split}  
\end{equation*}  

Remember the meaning of holding  $\gamma$, 
with 

\begin{equation*} 
\gamma=\gamma(\epsilon,\eta)
\end{equation*}  

then

\begin{equation*} 
d\gamma
=
\left(
\frac{\partial \gamma}{\partial \epsilon}
\right)_\eta
d\epsilon
+
\left(
\frac{\partial \gamma}{\partial \eta}
\right)_\epsilon
d\eta
=
0
\end{equation*}  

so  

\begin{equation*} 
d\gamma
=
\left(
\frac{\partial \gamma}{\partial \epsilon}
\right)_\eta
d\epsilon
+
\left(
\frac{\partial \gamma}{\partial \eta}
\right)_\epsilon
d\eta
=
0
\end{equation*}  

\begin{equation*} 
-
\dfrac{\left(
\frac{\partial \gamma}{\partial \epsilon}
\right)_\eta
}
{
\left(
\frac{\partial \gamma}{\partial \eta}
\right)_\epsilon
}
=
\frac{d\eta}{d\epsilon}
\bigg |_{\gamma=const}
=
\left(
\frac{\partial \eta}{\partial \epsilon}
\right)_\gamma
=
\frac{\gamma}{\sqrt{2\epsilon}}
,
\end{equation*} 

with 

\begin{equation*} 
\left(
\frac{\partial \gamma}{\partial \epsilon}
\right)_\eta
=
-
\frac{\gamma}{2\epsilon}
\qquad
\left(
\frac{\partial \gamma}{\partial \eta}
\right)_\epsilon
=\frac{1}{\sqrt{2\epsilon}}
\end{equation*} 

Finally, 

\begin{equation*} 
\left(
    \frac{\partial F(\chi,\gamma)}{\partial \chi}
\right)_\gamma
=
\left(
\frac{\partial F(\chi,\gamma)}{\partial \epsilon}
\right)_\eta
+
\left(
\frac{\partial F(\chi,\gamma)}{\partial \eta}
\right)_\epsilon
\frac{1}{\sqrt{2\epsilon}}
=
\left(
\frac{\partial F(\chi,\gamma)}{\partial \epsilon}
\right)_\eta
+
\frac{U}{\sqrt{2\epsilon}}
\end{equation*}  

\begin{equation*} 
\left(
    \frac{\partial F(\chi,\gamma)}{\partial \chi}
\right)_\gamma
=
\left(
\frac{\partial }{\partial \epsilon}
\right)_\eta
\left(
\frac{\psi}{\sqrt{2\epsilon}}
\right)
+
\frac{U}{\sqrt{2\epsilon}}
=
\frac{1}{\sqrt{2\epsilon}}
\left(
\frac{\partial \psi}{\partial \epsilon}
\right)_\eta
+
\frac{\psi}{{(2\epsilon)}^{3/2}}
+
\frac{U}{\sqrt{2\epsilon}}
\end{equation*}  

The similarity assumption request that 

\begin{equation*} 
F'=
\frac{\partial F(\chi,\gamma)}{\partial \gamma}
=
\frac{\partial F(\gamma) }{\partial \eta}
=
U
\end{equation*}  

It is an ansatz, and it only works if the governing equations
can be expressed in function only on -in this case- $\gamma$. 
Then,  

\begin{equation*} 
\frac{\partial \psi}{\partial \eta}
\left(
    \frac{\partial U}{\partial \chi}
    -
    \frac{\gamma}{2\epsilon}
    \frac{\partial U}{\partial \gamma}
\right)
-
\frac{\partial \psi}{\partial \epsilon}
\left(
    \frac{1}{\sqrt{2\epsilon}}
    \frac{\partial U}{\partial \gamma}
\right)
=
\frac{1}{\sqrt{2\epsilon}}
\frac{\partial}{\partial \gamma}
\left[
    \mu\rho
    \left(
    \frac{1}{\sqrt{2\epsilon}}
    \frac{\partial U}{\partial \gamma}
    \right)
\right]
,
\end{equation*} 

\begin{equation*} 
\frac{1}{\sqrt{2\epsilon}}
\frac{\partial }{\partial \gamma}
\left(
\sqrt{2\epsilon}
F
\right)
\left(
    \frac{\partial F'}{\partial \chi}
    -
    \frac{\gamma}{2\epsilon}
    F''
\right)
-
\frac{\partial }{\partial \epsilon}
\left(
        F \sqrt{2\epsilon}
     \right)
\left(
    \frac{1}{\sqrt{2\epsilon}}
    F''
     \right)
=
\frac{1}{\sqrt{2\epsilon}}
\frac{\partial}{\partial \gamma}
\left[
    \mu\rho
    \left(
    \frac{1}{\sqrt{2\epsilon}}
    F''
    \right)
\right]
,
\end{equation*} 

\begin{equation*} 
\frac{\partial }{\partial \gamma}
\left(
\sqrt{2\epsilon}
F
\right)
\left(
    \frac{\partial F'}{\partial \chi}
    -
    \frac{\gamma}{2\epsilon}
    F''
\right)
-
\left(
        \frac{\partial }{\partial \chi}
        \left(
        F \sqrt{2\epsilon}
        \right)
        -
        \frac{\gamma}{2\epsilon}
        \frac{\partial }{\partial \gamma}
        \left(
        F \sqrt{2\epsilon}
        \right)
\right)
F''
=
\frac{\partial}{\partial \gamma}
\left[
    \mu\rho
    \left(
    \frac{1}{\sqrt{2\epsilon}}
    F''
    \right)
\right]
,
\end{equation*} 

\begin{equation*} 
\left(
\sqrt{2\epsilon}
\frac{\partial F}{\partial \gamma}
\right)
\left(
    \frac{\partial F'}{\partial \chi}
    -
    \frac{\gamma}{2\epsilon}
    F''
\right)
-
\left(
        \sqrt{2\epsilon}
        \frac{\partial F }{\partial \chi}
        +
        F 
        \frac{1}{\sqrt{2\epsilon}}
        -
        \frac{\gamma}{\sqrt{2\epsilon}}
        \frac{\partial F}{\partial \gamma}
\right)
F''
=
\frac{\partial}{\partial \gamma}
\left[
    \mu\rho
    \left(
    \frac{1}{\sqrt{2\epsilon}}
    F''
    \right)
\right]
,
\end{equation*} 

\begin{equation*} 
\left(
2\epsilon
\frac{\partial F}{\partial \gamma}
\right)
\left(
    \frac{\partial F'}{\partial \chi}
    -
    \frac{\gamma}{2\epsilon}
    F''
\right)
-
\left(
        2\epsilon
        \frac{\partial F }{\partial \chi}
        +
        F 
        -
        \gamma
        \frac{\partial F}{\partial \gamma}
\right)
F''
=
\frac{\partial}{\partial \gamma}
\left[
    \mu\rho
    \left(
    F''
    \right)
\right]
,
\end{equation*} 

\begin{equation*} 
\left(
2\epsilon
\right)
\left(
    \frac{\partial F'}{\partial \chi}
    F'
        -
    \frac{\partial F }{\partial \chi}
    F''
\right)
-
\left(
        \cancel{
        \gamma
        \frac{\partial F}{\partial \gamma}
        }
        +
        F 
        -
        \cancel{
        \gamma
        \frac{\partial F}{\partial \gamma}
        }
\right)
F''
=
\frac{\partial}{\partial \gamma}
\left[
    \mu\rho
    \left(
    F''
    \right)
\right]
,
\end{equation*} 

\begin{equation*} 
2\epsilon
\left(
    F'
    \frac{\partial F'}{\partial \chi}
        -
    \frac{\partial F }{\partial \chi}
    F''
\right)
+
F
F''
+
\frac{\partial}{\partial \gamma}
\left[
    \mu\rho
    \left(
    F''
    \right)
\right]
=0
,
\end{equation*} 


By similarity ansatz

\begin{equation*} 
F
F''
+
\frac{\partial}{\partial \gamma}
\left[
    \mu\rho
    \left(
    F''
    \right)
\right]
=0
,
\end{equation*} 